\subsection*{5.2 $\tau$-sampling}
Toward establishing efficient decoding, we will need the notion of$\tau$-sampling.
\paragraph{Definition32 ($\tau$-sampling).} We say that a distribution$\mathbf W$ over$[n]^t$ is$\tau$-sampling if for any$i\in[t-1]$, $S\subseteq[n]$, and$X\subseteq[n]^i$,
\[\Cov_{\mathbf w\sim\mathbf W}[1[\mathbf w_{i+1}\in S],1[(\mathbf w_1,\ldots,\mathbf w_i)\in X]]\leq\tau.\]
To highlight the fact that it is indeed a (strong) sampling property, assume for simplicity that$\mathbf W$ is homogeneous. Fix any$i\in[t-1]$, $S\subseteq[n]$, and$X\subseteq[n]^i$. The property of$\tau$-sampling thus tells us we can use$\mathbf W$ to sample$S$ starting from any prefix. Namely, that
\[\left|\Pr_{\mathbf w\in\mathbf W}[\mathbf w_{i+1}\in S\mid(\mathbf w_1,\ldots,\mathbf w_i)\in X]-\rho(S)\right|\leq\frac\tau{\rho(X)}.\]
Note that the$t=2$ case corresponds to the setting of the expander mixing lemma.
A$\tau$-sampling distribution$\mathbf W$ satisfies the property of the ``Splittable Mixing Lemma'' of Jeronimo, Srivastava, and Tulsiani [20,Lemma4.6] for the family of ``$\pm1$ cut functions''. This fact is crucial us, and for completeness we establish this in AppendixA.
Given a$\tau$-sampling homogeneous$\mathbf W\sim[n]^t$, a standard hybrid argument shows the following.
\paragraph{Lemma33.} Let$\mathbf W\sim[n]^t$ be homogeneous and$\tau$-sampling, and let$z\in\mathbb F_2^n$ be arbitrary. Then,
\[\left|\mathbb E_{\mathbf w\in\mathbf W}[(-1)^{z_{\mathbf w_1}+\cdots+z_{\mathbf w_t}}]-\bias(z)^t\right|\leq2(t-1)\tau.\]
Thus, high-order mixing in particular implies that$\mathbf W$ is an$(\varepsilon_0,\varepsilon=\varepsilon_0^t+(t-1)\tau)$ parity sampler for all$\varepsilon_0$. However, for us, and also in[20], $t\cdot\tau$ is too large so we prove the parity sampling property separately.
We conclude our discussion about$\tau$-sampling by showing that our$\mathbf W$ is indeed$\tau$-sampling.
\paragraph{Lemma34.} For any$n,d,t\in\mathbb N$ and$\lambda>0$, if$H=(V=[n],E)$ is a$\lambda$-spectral$d$-regular3-uniform hypergraph, then$\mathbf W_{H,t}$ is$\tau=\lambda/4$-sampling.
\begin{proof}
As in Definition32, fix any$i\in[t-1]$, $S\subseteq[n]$ and$X\subseteq[n]^i$. Let$S_i(\mathbf w)$ be the indicator that$\mathbf w_{i+1}\in S$, and$X_i(\mathbf w)$ the indicator that$(\mathbf w_1,\ldots,\mathbf w_t)\in X$. Using the identity$\Cov[1-2a,1-2b]=4\Cov[a,b]$, it is sufficient to prove that
\[\Cov_{\mathbf w\sim\mathbf W_{H,t}}[(-1)^{S_i(\mathbf w)},(-1)^{X_i(\mathbf w)}]\leq\lambda.\]
Let$\mathbf v_0,\ldots,\mathbf v_t$, $\mathbf w_1,\ldots,\mathbf w_t$, and$\mathbf e_1,\ldots,\mathbf e_t$ be the random variables defined in the construction of$\mathbf W_{H,t}$. Furthermore, let$x,y,z\in\mathbb R^n$ be the vectors defined, for each$j\in[n]$, by
\[x_j\coloneqq\mathbb E[(-1)^{X_i(\mathbf w)}\mid\mathbf v_i=j],\qquad y_j\coloneqq(-1)^{1[j\in S]},\qquad z_j\coloneqq1/n.\]
Note that$\|z\|_2=1/\sqrt n$, $\|y\|_{\infty}=1$, and$\|x\|_2\leq\sqrt n$ (which follows from$\|x\|_{\infty}\leq1$). Our goal in this proof will be to show that the following equation holds:
\begin{multline*}
\Cov_{\mathbf w}[(-1)^{S_i(\mathbf w)},(-1)^{X_i(\mathbf w)}]\\
=\frac1d\sum_{(a,b,c)\in E}x_ay_bz_c-\frac1{n^2}\sum_{a\in[n]}x_a\sum_{a\in[n]}y_a\sum_{a\in[n]}z_a.\tag{8}
\end{multline*}
Once we do, the desired result follows from Definition12. In order to compute the covariance, we first expand:
\begin{align*}
&\mathbb E_{\mathbf w}[(-1)^{S_i(\mathbf w)}\cdot(-1)^{X_i(\mathbf w)}]\\
&=\frac1{nd}\sum_{e\in E}\mathbb E_{\mathbf w}[(-1)^{S_i(\mathbf w)}\cdot(-1)^{X_i(\mathbf w)}\mid\mathbf e_{i+1}=e]\qquad\text{(Claim24)}\\
&=\frac1{nd}\sum_{(a,b,c)\in E}(-1)^{1[b\in S]}\cdot\mathbb E_{\mathbf w}[(-1)^{X_i(\mathbf w)}\mid\mathbf v_i=a]\\
&=\frac1d\sum_{(a,b,c)\in E}x_ay_bz_c.\tag{9}
\end{align*}
Next, directly from the definition of$y$ and Claim24,
\[\mathbb E_{\mathbf w}[(-1)^{S_i(\mathbf w)}]=\frac1n\sum_{a\in[n]}y_a.\tag{10}\]
Similarly,
\[\mathbb E_{\mathbf w}[(-1)^{X_i(\mathbf w)}]=\frac1n\sum_{a\in[n]}x_a.\tag{11}\]
Equation(8) follows from Equations(9)–(11) and the fact that$\sum_{a\in[n]}z_a=1$. The desired result then follows Definition12.
\end{proof}
\subsection*{5.3 Decoding $C$}
We follow the work of Jernoimo et al. who gave a near-linear time list- and unique-decoding algorithm for Ta-Shma's code via an efficient weak regularity lemma, and prove that their algorithm also applies to our code as well. In the language of$\tau$-sampling distributions, they prove:\footnote{Jernoimo et al. prove their result under stronger requirements, however one can verify that the mixing requirement suffices.}
\paragraph{Theorem35 ([20]).} There exists a constant$c_{\mathrm{JST}}$ such that the following holds for any integers$d,t,k,n$ and any$\tau,\varepsilon_0,\varepsilon>0$. Let$C_0:\mathbb F_2^k\to\mathbb F_2^n$ be a code with bias at most$\varepsilon_0$ which is uniquely decodable to within distance$(1-\varepsilon_0)/4$ in time$T_0=T_0(n,\varepsilon_0)$. Let$\mathbf W\sim[n]^t$ be a homogeneous$\tau$-sampling distribution, let$C=\dsum_{\mathbf W}(C_0)$ be the corresponding direct sum lifting, and assume that the bias of$C$ is at most$\varepsilon$. Let$\beta$ be such that
\[\beta\geq\max\left\{\sqrt\varepsilon,\sqrt{c_{\mathrm{JST}}\cdot t^3\tau},2\left(\frac12+2\varepsilon_0\right)^t\right\}.\]
Then, there exists a randomized algorithm, which given$\tilde y\in\mathbb F_2^{|\mathbf W|}$, recovers the list$C\cap B(\tilde y,1/2-\beta)$ with probability at least$1-\frac1\varepsilon\cdot2^{-\Omega(\varepsilon_0^2n)}$ in time$\widetilde O(c_{\beta,t,\varepsilon_0}\cdot(|\mathbf W|+T_0))$, for$c_{\beta,t,\varepsilon_0}=(6/\varepsilon_0)^{2^{O(t^3/\beta^2)}}$. Moreover, if we are able to set$\beta=(1+\varepsilon)/4$, we can uniquely decode$C$ to within distance$(1-\varepsilon)/4$ with probability at least$1-2^{-\Omega(\varepsilon_0^2n)}$.
Plugging-in our code$C$ (using$\mathbf W_{H,t}$ and$C_0$ as described in Section5.1), we get our main result [Theorem36, stated above].
\begin{proof}[Proof of Theorem36]
By Lemma34, our parity sampler$\mathbf W$ which we use for$C$ is$\tau$-sampling for
\[\tau=\frac\lambda4=O\left(\frac{\log d}{\sqrt d}\right)=\widetilde O(2^{-2\sqrt{\log(1/\varepsilon)}}),\]
where we used the fact that our randomized construction of$H$ gives us$\lambda=O((1/\sqrt d)\log d)$ and the bound on$d$ from Equation(7). All that is left is to show how the two items follow from Theorem35. For the list decoding result, we take$\beta$ to be as small as possible. For us,
\[\left(\frac12+2\varepsilon_0\right)^t\leq2^{-\frac12\sqrt{\log(1/\varepsilon)}},\]
and$\sqrt{c_{\mathrm{JST}}\cdot t^3\tau}=\widetilde O(2^{-\sqrt{\log(1/\varepsilon)}})$. We can thus conclude that$\beta\leq2^{-\frac12\sqrt{\log(1/\varepsilon)}}$. (Again, we are assuming$\varepsilon$ is smaller than some small constant.) For the running time, note that
\[T_0=\exp(\operatorname{poly}(1/\varepsilon_0))n=\exp\left(2^{O(\sqrt{\log(1/\varepsilon)})}\right)\cdot k,\]
and$c_{\beta,t,\varepsilon_0}$ is thus triply-exponential in$\sqrt{\log(1/\varepsilon)}$, and overall$\widetilde O(c_{\beta,t,\varepsilon_0}\cdot(|\mathbf W|+T_0))=\widetilde O(c_{\beta,t,\varepsilon_0}\cdot k)$. For the unique decoding result, we take$\beta=(1+\varepsilon)/4$, and the running time becomes doubly-exponential in$\sqrt{\log(1/\varepsilon)}$.
\end{proof}
